\addtocontents{toc}{\protect\Needspace{3\baselineskip}}
\originalchapter{第三次测验复习：时间、曲率与视界}
下列五题组均来自 8.286 Fall 2013 第三次测验公开复习册，编号按原文件保留。第12题原文件自己说明超出当年课堂考核范围；这里把它作为有条件的化学势延伸，严格区分题设近似与真实费米分布。

\src{8.286 Quiz 3 Review Problem 8；官方复习解答。}
\begin{exercise}
从 $10^{-43},10^{-37},10^{-12},10^{-5},1$ 秒、4分钟、$10^4$--$10^6$年及若干十亿年选最合适时刻，分别给出：(a) 大爆炸核合成有关过程开始；(b) 结束；(c) $T\sim10^{16}\,\mathrm{GeV}$ 的大统一理论相变；(d) 原子复合；(e) $T\sim300\,\mathrm{MeV}$ 的夸克禁闭。原题是历史教学数量级选择，而非高精度热史。
\end{exercise}
\begin{solution}
按课程答案依次为 $1\,\mathrm s$、约4分钟、$10^{-37}\,\mathrm s$、$10^4$--$10^6$ 年、$10^{-5}\,\mathrm s$。推导主要靠辐射主导下 $T\propto t^{-1/2}$：把 $t\sim1\,\mathrm s$、$\kb T\sim1\,\mathrm{MeV}$ 作为课程锚点，温度升高因子 $f$ 对应时间缩短 $f^2$；$10^{16}\,\mathrm{GeV}=10^{19}\,\mathrm{MeV}$ 则粗略给 $10^{-38}\,\mathrm s$，备选项选 $10^{-37}$ 秒。复合需用氢结合、光子重子数比和膨胀史，不能仅凭 $13.6\,\mathrm{eV}$ 代入上述近似；课程在范围内约取二十万年。
\end{solution}

\src{8.286 Quiz 3 Review Problem 9(a)--(b), 原题为2004年 Quiz 3 Q3；官方复习解答。}
\begin{exercise}
在物质主导和辐射主导两个\emph{各自始终近乎平直}的理想宇宙中，从 Friedmann 式证明 $(\Omega-1)/\Omega$ 随宇宙时 $t$ 的幂律，并给出两种幂次。保持曲率参数的符号，说明近乎平直假设用在何处。
\end{exercise}
\begin{solution}
由 $\rho_c=3H^2/(8\pi G)$ 和 $H^2=8\pi G\rho/3-kc^2/a^2$，准确恒等式为
\[
\frac{\Omega-1}{\Omega}=\frac{3kc^2}{8\pi Ga^2\rho}.
\]
物质 $\rho\propto a^{-3}$，故比值 $\propto a$；近乎平直时 $a\propto t^{2/3}$，得 $\propto t^{2/3}$。辐射 $\rho\propto a^{-4}$，比值 $\propto a^2$；近乎平直时 $a\propto t^{1/2}$，得 $\propto t$。曲率符号由 $k$ 保留，方程说明偏离量的\emph{绝对值}随时间增长；若偏离已变大，不能继续使用近乎平直的 $a(t)$。
\end{solution}

\src{8.286 Quiz 3 Review Problem 10(a)--(e), 原题为2004年 Quiz 3 Q4；官方复习解答。}
\begin{exercise}
用原题给定的平直模型 $H_0=65\,\mathrm{km\,s^{-1}Mpc^{-1}}$、$\Omega_{m0}=0.35,\Omega_{\Lambda0}=0.65$，考察红移 $z=5.82$ 的 SDSS 类星体。检查历史说明中的五项数值：(a) 今天年龄约 $13.9$ 十亿年；(b) 所见光发出时宇宙约 $0.95$ 十亿年；(c) 今天标尺距离约 $27$ 十亿光年；(d) 发光时两者实际距离约 $4$ 十亿光年；(e) 今天由 Hubble 流定义的退行速度约 $1.8c$。可给定积分，无须计算数值。
\end{exercise}
\begin{solution}
令 $E(z)=\sqrt{\Omega_{m0}(1+z)^3+\Omega_{\Lambda0}}$。年龄与发射时刻分别为
\[
t_0=H_0^{-1}\int_0^\infty\frac{\dd z}{(1+z)E(z)},\qquad
t_e=H_0^{-1}\int_{5.82}^\infty\frac{\dd z}{(1+z)E(z)}.
\]
平直光路给今日实际距离 $D_0=(c/H_0)\int_0^{5.82}\dd z/E(z)$，发射时共动两端实际间距 $D_e=D_0/(1+5.82)$，今日退行速度定义为 $v_0=H_0D_0$。课程解答积分得到 $t_0\simeq13.88$ 十亿年、$t_e\simeq0.9509$ 十亿年、无量纲光路积分 $1.8099$、$D_0\simeq27.23$ \emph{十亿}光年、$D_e\simeq3.992$ 十亿光年、$v_0\simeq1.8099c$。原官方解答在一处把 $27.23$ 后写成“light-yr”而漏了十亿倍，量纲和其上下文均表明应是十亿光年。$v_0>c$ 是远隔共动点的 FLRW 距离变化率，不等于一个穿过本地实验室的物体速度。
\end{solution}

\src{8.286 Quiz 3 Review Problem 11(a)--(d), 原题为2007年 Quiz 3 Q3；官方复习解答。}
\begin{exercise}
一扰动模今天实际波长为 $\lambda_0$。题设其\emph{第二次} Hubble 长度穿越发生于 $30\,\mathrm s\ll t\ll50000$ 年的平直辐射主导期，当时辐射只有光子与三味中微子，$T_\nu/T_\gamma=(4/11)^{1/3}$。求(a) $c/H(t)$；(b) 写 $u=[g_1\pi^2/30](\kb T_\gamma)^4/(\hbar c)^3$ 时的 $g_1$；(c) $T_\gamma(t)$；(d) 由 $\lambda(t)=c/H(t)$ 求 $t_{H2}(\lambda_0)$，允许保留常数。原题历史背景使用 WMAP 3年参数 $H_0=73.5,\Omega_m=0.237,\Omega_\Lambda=0.763,T_{\gamma0}=2.725\,\mathrm K$，本解所需主要是今天的光子温度。
\end{exercise}
\begin{solution}
辐射期 $a\propto t^{1/2}$，故 $H=1/(2t)$、$c/H=2ct$。两偏振光子给 $2$，三味中微子及反中微子共六态给 $21/4$，但温度较低，所以
\[
g_1=2+\frac{21}{4}\left(\frac4{11}\right)^{4/3},\qquad
\kb T_\gamma(t)=\left(\frac{45\hbar^3c^5}{16\pi^3g_1G}\right)^{1/4}t^{-1/2}.
\]
题设这段之后不再有影响光子温度的热自由度湮灭，故 $aT_\gamma$ 守恒，$\lambda(t)=\lambda_0T_{\gamma0}/T_\gamma(t)$。把它与 $2ct$ 相等并平方，得到
\[
t_{H2}(\lambda_0)
=(\kb T_{\gamma0}\lambda_0)^2
\left(\frac{\pi^3g_1G}{45\hbar^3c^9}\right)^{1/2}.
\]
该式要求求得时刻确实在题定的辐射期；例如波长很大而跨越落在物质期时不能外推它。Hubble 半径的跨越也不是整个粒子视界的跨越。
\end{solution}

\src{8.286 Quiz 3 Review Problem 12(a)--(d), 原题为1998年 Quiz 4 Q4；官方复习册声明超出2013年考核范围。}
\begin{exercise}
早期热平衡的中子、质子静止能之差 $\Delta E=1.29\,\mathrm{MeV}$。允许化学势 $\mu_n\ne\mu_p$：(a) 修正 $n_n/n_p=\exp[-\Delta E/(\kb T)]$；(b) 用 $e^++n\leftrightarrow p+\bar\nu_e$、$\nu_e+n\leftrightarrow p+e^-$ 求 $\mu_n-\mu_p$ 与 $\mu_e,\mu_\nu$ 的关系；(c) 原题\emph{规定}中微子数密度按 $n_\nu=A\exp(\mu_\nu/\kb T)$、$n_{\bar\nu}=A\exp(-\mu_\nu/\kb T)$ 近似，其中 $A=3\zeta(3)(\kb T)^3/[4\pi^2(\hbar c)^3]$ 且 $\mu_e=0$，用 $n_{\bar\nu}/n_\nu$ 或两者差表示 $n_n/n_p$；(d) 反中微子过剩使氦份额怎样变化？
\end{exercise}
\begin{solution}
非相对论 Maxwell--Boltzmann 核子气体直接给
\[
\frac{n_n}{n_p}=\left(\frac{m_n}{m_p}\right)^{3/2}
\exp\!\left[\frac{\mu_n-\mu_p-\Delta E}{\kb T}\right].
\]
反粒子化学势相反，平衡条件 $-\mu_e+\mu_n=\mu_p-\mu_\nu$，即 $\mu_n-\mu_p=\mu_e-\mu_\nu$。略去质量比的千分量并令 $\mu_e=0$，则
\[
\frac{n_n}{n_p}=\exp[-\Delta E/(\kb T)-\mu_\nu/(\kb T)]
=\sqrt{\frac{n_{\bar\nu}}{n_\nu}}\,
\exp[-\Delta E/(\kb T)].
\]
若给差值 $\Delta n=n_{\bar\nu}-n_\nu$，令 $x=\exp[-\mu_\nu/(\kb T)]>0$，由 $\Delta n=A(x-x^{-1})$ 得
$x=[\Delta n+\sqrt{\Delta n^2+4A^2}]/(2A)$，代前式即可。反中微子过剩使 $x>1$、$n/p$ 上升，在其余核反应条件固定时更多中子能组成 $^4\mathrm{He}$。
原题把有化学势的相对论中微子数密度简化为 $A\ee^{\mu/\kb T}$；精确 Fermi--Dirac 积分一般不等于该指数式。本问按题设使用它，不能把平方根关系推广为任意退化中微子气体的精确公式。
\end{solution}
